% ============================================================================
% CAPÍTULO 6: RESULTADOS
% ============================================================================
%
% REQUISITOS SEGÚN LA FACULTAD:
%
% Este apartado debe presentar:
%
% 1. RESULTADOS DEL PROCESO DE VALIDACIÓN DEL INSTRUMENTO:
%    • Validación por juicio de expertos
%    • Confiabilidad y consistencia interna
%    • Otras técnicas de validación empleadas
%
% 2. CARACTERIZACIÓN DE LA POBLACIÓN Y/O UNIDADES DE ANÁLISIS
%
% 3. RESULTADOS ACORDES AL PLAN DE ANÁLISIS:
%    • Análisis univariado (frecuencias y análisis de medias)
%    • Análisis bivariado (comparación, correlación, regresión lineal)
%    • Análisis multivariado (regresión múltiple y AFE)
%
% 4. INCLUIR UN CUADRO O FIGURA POR PÁGINA ACORDE AL FORMATO APA
%
% 5. IDENTIFICAR LOS HALLAZGOS MÁS IMPORTANTES
%
% 6. NO REPETIR LOS RESULTADOS EN EL TEXTO SI ESTÁN EN UN CUADRO O GRÁFICA
%
% 7. TESIS: En un párrafo se describe la principal aportación de la investigación
%    derivado de la constatación de la hipótesis + un diagrama que explique la tesis
%
% ============================================================================

\chapter{Resultados}
\label{cap:resultados}

% INSTRUCCIONES GENERALES:
% - Presenta resultados OBJETIVOS (sin interpretación)
% - La interpretación va en DISCUSIÓN
% - Usa tablas y figuras con formato APA
% - NO repitas en texto lo que está en tablas/figuras
% - Identifica y resalta hallazgos importantes
% - Al final: presenta tu TESIS con diagrama

En este capítulo se presentan los resultados obtenidos de acuerdo con los objetivos planteados y el plan de análisis descrito en la metodología.

% ----------------------------------------------------------------------------
% SECCIÓN 6.1: Validación de Instrumentos
% ----------------------------------------------------------------------------
\section{Resultados de Validación de Instrumentos}
\label{sec:validacion_resultados}

% SOLO SI APLICA: Si validaste o adaptaste instrumentos

\subsection{Validación de Contenido (Juicio de Expertos)}
\label{subsec:resultados_juicio_expertos}

[Si realizaste validación por expertos, presenta los resultados]

Se sometió el instrumento a evaluación de \textit{[n]} expertos. La \Cref{tab:validacion_expertos} muestra los resultados del Índice de Validez de Contenido (IVC) por ítem.

\begin{table}[htbp]
    \centering
    \caption{Índice de Validez de Contenido por ítem}
    \label{tab:validacion_expertos}
    \small
    \begin{tabular}{@{}clcccc@{}}
        \toprule
        \textbf{Ítem} & \textbf{Enunciado} & \textbf{Claridad} & \textbf{Pertinencia} & \textbf{Suficiencia} & \textbf{IVC} \\
        \midrule
        1 & \textit{[Enunciado resumido]} & 0.90 & 0.95 & 0.85 & 0.90 \\
        2 & \textit{[...]} & ... & ... & ... & ... \\
        \midrule
        \multicolumn{5}{l}{\textbf{IVC Global del Instrumento}} & \textbf{[valor]} \\
        \bottomrule
    \end{tabular}
\end{table}

El IVC global del instrumento fue de \textit{[valor]}, considerado \textit{[excelente/aceptable según criterio > 0.80]}.

\subsection{Análisis de Confiabilidad}
\label{subsec:resultados_confiabilidad}

[Presenta resultados de confiabilidad]

\textbf{Consistencia Interna (Alfa de Cronbach):}

La \Cref{tab:alfa_cronbach} presenta los valores de confiabilidad del instrumento.

\begin{table}[htbp]
    \centering
    \caption{Confiabilidad del instrumento (Alfa de Cronbach)}
    \label{tab:alfa_cronbach}
    \begin{tabular}{@{}lcc@{}}
        \toprule
        \textbf{Dimensión} & \textbf{N° Ítems} & \textbf{Alfa de Cronbach} \\
        \midrule
        Dimensión 1 & [n] & [valor] \\
        Dimensión 2 & [n] & [valor] \\
        Dimensión 3 & [n] & [valor] \\
        \midrule
        \textbf{Escala Global} & \textbf{[total]} & \textbf{[valor]} \\
        \bottomrule
    \end{tabular}
\end{table}

Todos los valores de $\alpha$ fueron superiores a 0.70, indicando consistencia interna aceptable.

\subsection{Validación de Constructo (si aplica)}
\label{subsec:resultados_constructo}

[Si hiciste AFE con muestra amplia]

Se realizó Análisis Factorial Exploratorio (AFE). La prueba de Kaiser-Meyer-Olkin (KMO = \textit{[valor]}) y la prueba de esfericidad de Bartlett ($\chi^2$ = \textit{[valor]}, $p < 0.001$) indicaron adecuación de los datos para el análisis factorial.

Se extrajeron \textit{[n]} factores que explicaron el \textit{[\%]} de la varianza total.

% ----------------------------------------------------------------------------
% SECCIÓN 6.2: Caracterización de la Población
% ----------------------------------------------------------------------------
\section{Caracterización de la Población de Estudio}
\label{sec:caracterizacion_poblacion}

\subsection{Características Sociodemográficas}
\label{subsec:caracteristicas_sociodemograficas}

La muestra estuvo conformada por \textit{n = [número]} participantes. La \Cref{tab:sociodemograficos} presenta las características sociodemográficas.

\begin{table}[htbp]
    \centering
    \caption{Características sociodemográficas de la población de estudio (N=[n])}
    \label{tab:sociodemograficos}
    \begin{tabular}{@{}lcc@{}}
        \toprule
        \textbf{Variable} & \textbf{n} & \textbf{\%} \\
        \midrule
        \textbf{Sexo} & & \\
        \quad Hombre & [n] & [\%] \\
        \quad Mujer & [n] & [\%] \\
        \midrule
        \textbf{Edad (años)} & & \\
        \quad 18-29 & [n] & [\%] \\
        \quad 30-39 & [n] & [\%] \\
        \quad 40-49 & [n] & [\%] \\
        \quad $\geq$50 & [n] & [\%] \\
        \quad Media (DE) & \multicolumn{2}{c}{[media] $\pm$ [DE]} \\
        \midrule
        \textbf{Escolaridad} & & \\
        \quad Primaria & [n] & [\%] \\
        \quad Secundaria & [n] & [\%] \\
        \quad Preparatoria & [n] & [\%] \\
        \quad Licenciatura o más & [n] & [\%] \\
        \midrule
        \textbf{[Otra variable]} & & \\
        \quad [Categoría 1] & [n] & [\%] \\
        \quad [Categoría 2] & [n] & [\%] \\
        \bottomrule
    \end{tabular}
    \\[5pt]
    \footnotesize{\textit{Nota:} DE = Desviación Estándar}
\end{table}

La mayoría de los participantes fueron \textit{[describe brevemente el hallazgo más relevante sin repetir todos los números]}.

% ----------------------------------------------------------------------------
% SECCIÓN 6.3: Análisis Univariado
% ----------------------------------------------------------------------------
\section{Análisis Descriptivo de Variables (Análisis Univariado)}
\label{sec:analisis_univariado}

\subsection{Variables Cuantitativas}
\label{subsec:variables_cuantitativas}

La \Cref{tab:descriptivos_cuantitativos} presenta las medidas de tendencia central y dispersión de las variables cuantitativas.

\begin{table}[htbp]
    \centering
    \caption{Estadísticos descriptivos de variables cuantitativas}
    \label{tab:descriptivos_cuantitativos}
    \begin{tabular}{@{}lccccc@{}}
        \toprule
        \textbf{Variable} & \textbf{Media} & \textbf{DE} & \textbf{Mediana} & \textbf{Mín} & \textbf{Máx} \\
        \midrule
        [Variable 1] & [valor] & [valor] & [valor] & [valor] & [valor] \\
        [Variable 2] & [valor] & [valor] & [valor] & [valor] & [valor] \\
        [Variable 3] & [valor] & [valor] & [valor] & [valor] & [valor] \\
        \bottomrule
    \end{tabular}
    \\[5pt]
    \footnotesize{\textit{Nota:} DE = Desviación Estándar; Mín = Mínimo; Máx = Máximo}
\end{table}

[Describe brevemente los hallazgos más relevantes]

\subsection{Variables Cualitativas}
\label{subsec:variables_cualitativas}

[Presenta frecuencias y porcentajes de tus variables principales]

La \Cref{fig:distribucion_variable1} muestra la distribución de \textit{[variable principal]}.

\begin{figure}[htbp]
    \centering
    \includegraphics[width=0.7\textwidth]{figuras/distribucion_variable1.png}
    \caption{Distribución de [nombre de la variable]}
    \label{fig:distribucion_variable1}
\end{figure}

Se observó que \textit{[hallazgo principal sin repetir los números exactos de la gráfica]}.

% ----------------------------------------------------------------------------
% SECCIÓN 6.4: Análisis Bivariado
% ----------------------------------------------------------------------------
\section{Análisis de Relaciones entre Variables (Análisis Bivariado)}
\label{sec:analisis_bivariado}

\subsection{Relación entre Variable Dependiente e Independientes}
\label{subsec:relacion_variables_principales}

[Aquí presentas las relaciones entre tus variables principales]

\textbf{[Variable Independiente 1] vs [Variable Dependiente]:}

La \Cref{tab:relacion_vi1_vd} muestra la relación entre \textit{[VI1]} y \textit{[VD]}.

\begin{table}[htbp]
    \centering
    \caption{Relación entre [VI1] y [VD]}
    \label{tab:relacion_vi1_vd}
    \begin{tabular}{@{}lccccc@{}}
        \toprule
        \textbf{[VI1]} & \multicolumn{2}{c}{\textbf{[VD Categoría 1]}} & \multicolumn{2}{c}{\textbf{[VD Categoría 2]}} & \textbf{Valor p} \\
        \cmidrule(lr){2-3} \cmidrule(lr){4-5}
         & n & \% & n & \% & \\
        \midrule
        [Categoría A] & [n] & [\%] & [n] & [\%] & \multirow{3}{*}{[valor p]} \\
        [Categoría B] & [n] & [\%] & [n] & [\%] & \\
        [Categoría C] & [n] & [\%] & [n] & [\%] & \\
        \bottomrule
    \end{tabular}
    \\[5pt]
    \footnotesize{\textit{Nota:} Prueba de $\chi^2$. Significancia estadística: $p < 0.05$}
\end{table}

Se encontró una relación \textit{[estadísticamente significativa/no significativa]} entre \textit{[VI1]} y \textit{[VD]} ($\chi^2$ = \textit{[valor]}, $p$ = \textit{[valor]}).

\subsection{Correlaciones entre Variables Cuantitativas}
\label{subsec:correlaciones}

[Si tienes correlaciones]

La \Cref{tab:matriz_correlaciones} presenta la matriz de correlaciones entre las variables cuantitativas.

\begin{table}[htbp]
    \centering
    \caption{Matriz de correlaciones (Coeficiente de Pearson/Spearman)}
    \label{tab:matriz_correlaciones}
    \begin{tabular}{@{}lcccc@{}}
        \toprule
        \textbf{Variable} & \textbf{Variable 1} & \textbf{Variable 2} & \textbf{Variable 3} & \textbf{Variable 4} \\
        \midrule
        Variable 1 & 1.00 & & & \\
        Variable 2 & [r]* & 1.00 & & \\
        Variable 3 & [r] & [r]** & 1.00 & \\
        Variable 4 & [r]* & [r] & [r]** & 1.00 \\
        \bottomrule
    \end{tabular}
    \\[5pt]
    \footnotesize{* $p < 0.05$; ** $p < 0.01$}
\end{table}

Se encontraron correlaciones significativas entre \textit{[menciona las más relevantes]}.

% ----------------------------------------------------------------------------
% SECCIÓN 6.5: Análisis Multivariado
% ----------------------------------------------------------------------------
\section{Análisis Multivariado}
\label{sec:analisis_multivariado}

[Presenta tus modelos de regresión o análisis factorial]

\subsection{Modelo de Regresión [Lineal/Logística]}
\label{subsec:regresion}

Para identificar los predictores de \textit{[variable dependiente]}, se realizó un análisis de regresión \textit{[lineal múltiple/logística]}.

La \Cref{tab:regresion} presenta los resultados del modelo final.

\begin{table}[htbp]
    \centering
    \caption{Modelo de regresión [tipo] para predecir [variable dependiente]}
    \label{tab:regresion}
    \begin{tabular}{@{}lcccc@{}}
        \toprule
        \textbf{Variable} & \textbf{Coeficiente B} & \textbf{Error Estándar} & \textbf{IC 95\%} & \textbf{Valor p} \\
        \midrule
        (Constante) & [valor] & [valor] & [[inf] - [sup]] & [p] \\
        [Predictor 1] & [valor] & [valor] & [[inf] - [sup]] & [p]* \\
        [Predictor 2] & [valor] & [valor] & [[inf] - [sup]] & [p]** \\
        [Predictor 3] & [valor] & [valor] & [[inf] - [sup]] & [p] \\
        \midrule
        \multicolumn{5}{l}{$R^2$ ajustado = [valor] / $F$ = [valor], $p < 0.001$} \\
        \bottomrule
    \end{tabular}
    \\[5pt]
    \footnotesize{* $p < 0.05$; ** $p < 0.01$. IC = Intervalo de Confianza}
\end{table}

El modelo explicó el \textit{[\%]} de la varianza de \textit{[variable dependiente]} ($R^2$ ajustado = \textit{[valor]}). Los predictores significativos fueron \textit{[lista]}.

% ----------------------------------------------------------------------------
% SECCIÓN 6.6: Resultados por Objetivo Específico
% ----------------------------------------------------------------------------
\section{Resultados por Objetivo Específico}
\label{sec:resultados_por_objetivo}

[Organiza tus resultados según tus objetivos específicos]

\subsection{Objetivo Específico 1: \textit{[Enunciado del OE1]}}
\label{subsec:resultados_oe1}

[Presenta los resultados que responden a este objetivo]

\subsection{Objetivo Específico 2: \textit{[Enunciado del OE2]}}
\label{subsec:resultados_oe2}

[Presenta los resultados que responden a este objetivo]

[Continúa con todos tus objetivos específicos]

% ----------------------------------------------------------------------------
% SECCIÓN 6.7: Síntesis de Hallazgos Principales
% ----------------------------------------------------------------------------
\section{Síntesis de Hallazgos Principales}
\label{sec:sintesis_hallazgos}

Los hallazgos más relevantes de esta investigación fueron:

\begin{enumerate}
    \item \textbf{Hallazgo 1:} \textit{[Descripción concisa]}
    \item \textbf{Hallazgo 2:} \textit{[Descripción concisa]}
    \item \textbf{Hallazgo 3:} \textit{[Descripción concisa]}
    \item \textbf{Hallazgo 4:} \textit{[Descripción concisa]}
\end{enumerate}

% ----------------------------------------------------------------------------
% SECCIÓN 6.8: TESIS (Obligatoria según la facultad)
% ----------------------------------------------------------------------------
\section{Tesis}
\label{sec:tesis}

% INSTRUCCIONES:
% - Describe en UN PÁRRAFO la principal aportación de tu investigación
% - Debe derivarse de la constatación de tu hipótesis
% - Incluye un DIAGRAMA que explique visualmente tu tesis

Con base en los resultados obtenidos y la constatación de la hipótesis de investigación, se establece la siguiente tesis:

\begin{center}
\fbox{%
\parbox{0.9\textwidth}{%
\textbf{TESIS:}

\textit{[Escribe aquí en un párrafo la principal aportación de tu investigación. Debe ser la conclusión central que se deriva de tus resultados y que responde a tu hipótesis]}

\textbf{Ejemplo:}

\textit{La implementación de [intervención/modelo propuesto] resulta en una mejora significativa de [variable de resultado] en [población específica], mediada por [mecanismo identificado]. Esta relación se mantiene incluso controlando por [variables confusoras], lo cual demuestra que [conclusión teórica o práctica principal].}
}%
}
\end{center}

\subsection{Diagrama Explicativo de la Tesis}
\label{subsec:diagrama_tesis}

La \Cref{fig:diagrama_tesis} presenta un modelo visual que sintetiza la tesis de esta investigación.

\begin{figure}[htbp]
    \centering
    \includegraphics[width=0.9\textwidth]{figuras/diagrama_tesis.png}
    \caption{Modelo explicativo de la tesis: relación entre [variables principales]}
    \label{fig:diagrama_tesis}
\end{figure}

[Explica brevemente el diagrama]

El diagrama muestra que \textit{[explica las flechas, relaciones, mediadores, moderadores que aparecen en tu diagrama]}.

% ============================================================================
% NOTAS FINALES:
% ============================================================================
%
% RECUERDA:
% 1. Resultados OBJETIVOS (sin interpretación)
% 2. NO repetir en texto lo que está en tablas/figuras
% 3. Resaltar hallazgos IMPORTANTES
% 4. Incluir la TESIS con su diagrama (obligatorio)
% 5. Todas las tablas/figuras con formato APA
% 6. Organizar por objetivos específicos
% 7. Incluir todos los niveles de análisis (univariado, bivariado, multivariado)
%
% ============================================================================

